\chapter{مطالعات تطبیقی و ارزیابی شاخص‌ها}
\label{app:comparison}

در این بخش، طرح مهستان در کنار تجربه‌های موفق و شکست‌خورده‌ی جهانی قرار می‌گیرد تا نقاط قوت و ریسک‌های آن شفاف شود.

\begin{landscape}
\section{جدول مقایسه‌ای جامع گذارهای دموکراتیک}

\begin{table}[htbp]
\centering
\caption{جدول مقایسه‌ای جامع گذارهای دموکراتیک}
\small
\begin{tabular}{R{2.2cm} C{1.1cm} C{1.1cm} C{1.1cm} C{1.1cm} C{1.1cm} C{1.1cm} C{1.1cm} C{1.1cm} C{1.1cm}}
\hline
\textbf{شاخص} & \textbf{مهستان} & \textbf{تونس} & \textbf{آفریقای جنوبی} & \textbf{اسپانیا} & \textbf{لهستان} & \textbf{شیلی} & \textbf{لیبی} & \textbf{مصر} & \textbf{عراق} \\
\hline
مشروعیت دموکراتیک & ۹ & ۸ & ۸ & ۷ & ۵ & ۷ & ۳ & ۴ & ۲ \\
عدالت انتقالی & ۸ & ۶ & ۹ & ۲ & ۴ & ۵ & ۱ & ۲ & ۲ \\
اصلاح امنیتی (SSR) & ۸ & ۴ & ۶ & ۵ & ۴ & ۳ & ۱ & ۲ & ۱ \\
حقوق اقلیت‌ها & ۹ & ۵ & ۹ & ۸ & ۵ & ۴ & ۱ & ۳ & ۳ \\
قانون اساسی‌گرایی & ۱۰ & ۶ & ۹ & ۸ & ۶ & ۵ & ۲ & ۴ & ۴ \\
مدیریت اقتصادی & ۸ & ۴ & ۵ & ۸ & ۵ & ۶ & ۲ & ۳ & ۲ \\
دموکراسی مستقیم & ۹ & ۴ & ۵ & ۷ & ۳ & ۷ & ۱ & ۲ & ۱ \\
\hline
\textbf{میانگین وزنی} & \textbf{۸.۷} & \textbf{۵.۹} & \textbf{۷.۴} & \textbf{۶.۳} & \textbf{۵.۱} & \textbf{۵.۲} & \textbf{۱.۸} & \textbf{۲.۹} & \textbf{۲.۳} \\
\hline
\end{tabular}
\end{table}
\end{landscape}

\section{توزیع کیفی امتیازات طرح مهستان}

نمودار زیر نشان‌دهنده‌ی توازن میان ابعاد مختلف طرح و تمرکز ویژه بر قانون اساسی‌گرایی و مشروعیت است.

\begin{center}
\begin{tikzpicture}
    \pie[rotate=0, radius=3, text=inside, color={MainBlue!80, LightBlue!80, SoftGreen!80, SoftPurple!80, SoftGold!80, SlateGray!80, AccentGold!80}]
    {15/مشروعیت, 12/فراگیری, 12/قانون اساسی, 10/عدالت, 10/امنیت, 10/اقتصاد, 31/سایر عوامل}
\end{tikzpicture}
\end{center}

\begin{modernbox}{تحلیل عدم موفقیت در عراق و لیبی}
مطالعه‌ی این دو کشور نشان می‌دهد که «خلأ نهادی» و «مداخله‌ی بی‌رویه‌ی خارجی» بزرگ‌ترین قاتل هر گذار است. طرح مهستان با تکیه بر «نمایندگی سریع ملی» و «حفاظت از زیرساخت‌های اداری» راه میانه را برگزیده است.
\end{modernbox}
